\Qhead{1.1}{Линейные колебательные системы. Силовое и параметрическое воздействие}{Линейные колебательные системы с одной степенью свободы. Силовое и параметрическое воздействие на линейные и слабонелинейные колебательные системы.}
\Subfile{1.1.1}{sections/01_oscillations/q01/s01}
\Subfile{1.1.2}{sections/01_oscillations/q01/s02}
\Subfile{1.1.3}{sections/01_oscillations/q01/s03}

\recap{Линейная система: одна координата, энергия и потери. Силовое воздействие даёт АЧХ и фазу; рост ограничен затуханием. Слабая нелинейность изгибает кривую. Параметрическое воздействие меняет коэффициент, работает около двойной частоты и имеет порог $2/Q$.}
\board{Осциллятор; резонансный «купол» и фаза; наклон Дуффинга; язык параметрической неустойчивости.}
\oral{Почему силовой отклик существует без порога, а параметрическая генерация требует превысить потери?}
\bait{Можно оставить два фазовых состояния параметрического колебания: сдвиг на $\pi$ сохраняет его отношение к накачке на двойной частоте.}
\sources{Пятаков~\cite{pyatakov}, PDF-с.~27--31, 43--52: вынужденный и параметрический резонансы. Изложение переработано по вопросу~1 конспекта коллеги; нормировки и энергетические равенства проверены из выписанных моделей.}
